% Fichier généré : ne pas modifier à la main.
\documentclass[11pt,a4paper]{article}
\usepackage[T1]{fontenc}
\usepackage[utf8]{inputenc}
\usepackage[french]{babel}
\usepackage[margin=19mm]{geometry}
\usepackage{lmodern}
\usepackage{microtype}
\usepackage[table]{xcolor}
\usepackage{array,longtable,booktabs}
\definecolor{ink}{HTML}{163F59}
\definecolor{soft}{HTML}{EFF4F7}
\renewcommand{\familydefault}{\sfdefault}
\setlength{\parindent}{0pt}
\setlength{\parskip}{5pt}
\setlength{\tabcolsep}{6pt}
\renewcommand{\arraystretch}{1.18}
\pagestyle{plain}
\begin{document}
{\color{ink}\Huge\bfseries Revoir une notion}\par
{\Large Capsules vidéo et repères dans les livres}\par


Manuel et cahier MAT0339 : édition du 2026-10-10.\par

Après le cours, repérez la notion à revoir, puis revenez aux sections du
manuel et aux exercices du cahier. Le numéro à gauche identifie la vidéo.
\textbf{Les renvois du manuel et du cahier se lisent séparément.}

\S\ désigne une section; « Ch. » désigne un chapitre lorsque le rapprochement
est seulement thématique. Chaque vidéo soutient une partie des sections citées :
elle ne remplace ni leur étude complète ni la recherche des exercices.
Ce répertoire n'atteste pas que toutes les vidéos sont disponibles ou validées.


\section*{Notions du cours}


\small
\begin{longtable}{@{}p{0.055\textwidth}p{0.525\textwidth}p{0.17\textwidth}p{0.17\textwidth}@{}}
\rowcolor{soft}\textbf{No} & \textbf{Vidéo et point de révision} & \textbf{Manuel} & \textbf{Cahier} \\
\midrule
\endfirsthead
\rowcolor{soft}\textbf{No} & \textbf{Vidéo et point de révision} & \textbf{Manuel} & \textbf{Cahier} \\
\midrule
\endhead
\midrule
\multicolumn{4}{r}{\footnotesize Suite à la page suivante}\\
\endfoot
\bottomrule
\endlastfoot

01 & \textbf{Variables et polynômes}\newline {\footnotesize Variables, coefficients et structure des polynômes.} & \S\ 2.1 & \S\ 2.1 \\
\addlinespace[7pt]

02 & \textbf{Inéquations et nombres négatifs}\newline {\footnotesize Comprendre pourquoi un facteur négatif inverse le sens de l’inégalité.} & \S\ 2.4 & \S\ 2.4 \\
\addlinespace[7pt]

03 & \textbf{Racine carrée et valeur absolue}\newline {\footnotesize Distinguer une racine carrée, les solutions d’une équation et la valeur absolue.} & \S\S\ 1.4, 1.6 & \S\ 1.4 \\
\addlinespace[7pt]

04 & \textbf{Synthèse des opérations sur les nombres réels}\newline {\footnotesize Priorités, signes, parenthèses et calcul sur les fractions.} & \S\S\ 1.2, 1.3 & \S\S\ 1.2, 1.3 \\
\addlinespace[7pt]

05 & \textbf{Domaine et image d'une fonction}\newline {\footnotesize Distinguer les entrées permises, le codomaine et les sorties atteintes.} & \S\ 3.1 & \S\ 3.1 \\
\addlinespace[7pt]

06 & \textbf{Modèles affines et quadratiques}\newline {\footnotesize Reconnaître une variation affine ou quadratique.} & \S\S\ 4.1, 4.2, 4.5 & \S\S\ 4.1, 4.2, 4.5 \\
\addlinespace[7pt]

07 & \textbf{Lire les propriétés d'un graphe}\newline {\footnotesize Lire croissance, extrema et symétries sur un graphe.} & \S\S\ 3.2, 3.7 & \S\S\ 3.2, 3.7 \\
\addlinespace[7pt]

08 & \textbf{Fonction définie par morceaux}\newline {\footnotesize Choisir une règle selon l’intervalle auquel appartient l’entrée.} & \S\ 3.8 & \S\ 3.8 \\
\addlinespace[7pt]

09 & \textbf{Opérations sur les fonctions}\newline {\footnotesize Accent sur l’addition et la soustraction graphiques; la section étudie aussi d’autres opérations.} & \S\ 3.4 & \S\ 3.4 \\
\addlinespace[7pt]

10 & \textbf{Composition de fonctions}\newline {\footnotesize Faire passer une sortie dans une seconde fonction et vérifier son domaine.} & \S\ 3.4 & \S\ 3.4 \\
\addlinespace[7pt]

11 & \textbf{Fonction réciproque}\newline {\footnotesize Retrouver l’entrée et comprendre la réflexion du graphe.} & \S\ 3.5 & \S\ 3.5 \\
\addlinespace[7pt]

12 & \textbf{Racines et hauteur zéro}\newline {\footnotesize Relier une racine à une intersection avec l’axe horizontal; la section traite les quadratiques.} & \S\ 4.3 & \S\ 4.3 \\
\addlinespace[7pt]

13 & \textbf{Multiplicité des racines}\newline {\footnotesize Expliquer pourquoi la courbe traverse ou touche l’axe.} & \S\S\ 5.3, 5.4 & \S\ 5.3 \\
\addlinespace[7pt]

14 & \textbf{Compléter le carré}\newline {\footnotesize Passer à la forme canonique et comprendre la formule quadratique.} & \S\S\ 4.2, 4.3 & \S\S\ 4.2, 4.3 \\
\addlinespace[7pt]

15 & \textbf{Fonctions rationnelles et leurs asymptotes}\newline {\footnotesize Conserver le domaine et distinguer un trou des asymptotes.} & \S\S\ 5.2, 5.3, 5.4 & \S\S\ 5.1, 5.2, 5.3 \\
\addlinespace[7pt]

16 & \textbf{Exponentielles et logarithmes}\newline {\footnotesize Relier exponentielle, logarithme et propriétés de calcul.} & \S\S\ 6.1, 6.2, 6.3 & \S\S\ 6.1, 6.2, 6.3 \\
\addlinespace[7pt]

17 & \textbf{Principe fondamental du dénombrement}\newline {\footnotesize Distinguer les étapes successives des cas disjoints.} & \S\ 7.1 & \S\ 7.1 \\
\addlinespace[7pt]

18 & \textbf{Permutation, arrangement et combinaison}\newline {\footnotesize Choisir une méthode de comptage selon les rôles et l’ordre.} & \S\S\ 7.2, 7.3, 7.4 & \S\S\ 7.2, 7.3, 7.4 \\
\addlinespace[7pt]

19 & \textbf{Répétitions en dénombrement}\newline {\footnotesize Distinguer les choix réutilisables des objets identiques à réordonner.} & \S\S\ 7.1, 7.2, 7.3 & \S\S\ 7.1, 7.2, 7.3 \\
\addlinespace[7pt]

20 & \textbf{Modèles probabilistes élémentaires}\newline {\footnotesize Construire un univers et vérifier quand les issues sont équiprobables.} & \S\S\ 8.1, 8.2 & \S\S\ 8.1, 8.2 \\
\addlinespace[7pt]

21 & \textbf{Déplacement et composantes}\newline {\footnotesize Relier un déplacement, ses composantes et sa longueur.} & \S\S\ 9.1, 9.3 & \S\S\ 9.1, 9.3 \\
\addlinespace[7pt]

22 & \textbf{Opérations sur les vecteurs}\newline {\footnotesize Additionner, soustraire et multiplier un vecteur par un scalaire.} & \S\ 9.2 & \S\ 9.2 \\
\addlinespace[7pt]

23 & \textbf{Vecteurs dans R\textsuperscript{2} et R\textsuperscript{3}}\newline {\footnotesize Calculer un déplacement, une distance et un milieu.} & \S\S\ 9.3, 9.8 & \S\S\ 9.3, 9.8 \\
\addlinespace[7pt]

24 & \textbf{Équations de la droite et du plan}\newline {\footnotesize Passer d’une description géométrique à un test d’appartenance.} & \S\S\ 10.1, 10.4 & \S\S\ 10.1, 10.4 \\
\addlinespace[7pt]

25 & \textbf{Lire et appliquer une matrice}\newline {\footnotesize Lire les dimensions et appliquer une matrice à un vecteur.} & \S\S\ 11.1, 11.5 & \S\S\ 11.1, 11.5 \\
\addlinespace[7pt]

26 & \textbf{Opérations sur les matrices}\newline {\footnotesize Contrôler les dimensions et l’ordre d’un produit.} & \S\S\ 11.2, 11.3 & \S\S\ 11.2, 11.3 \\
\addlinespace[7pt]

27 & \textbf{Déterminant et matrice inverse}\newline {\footnotesize Relier déterminant, aire, orientation et possibilité d’inverser.} & \S\S\ 11.7, 12.5 & \S\S\ 11.7, 12.5 \\
\addlinespace[7pt]

28 & \textbf{Méthode d’élimination de variables}\newline {\footnotesize Éliminer une inconnue tout en conservant un système équivalent.} & \S\S\ 12.2, 12.8 & \S\S\ 12.2, 12.8 \\
\addlinespace[7pt]

29 & \textbf{Règle de Cramer et systèmes par déterminants}\newline {\footnotesize Utiliser les déterminants lorsque celui du système est non nul.} & \S\ 12.6 & \S\ 12.6 \\
\addlinespace[7pt]

30 & \textbf{Programmation linéaire à deux variables}\newline {\footnotesize Distinguer les points admissibles de ceux qui optimisent l’objectif.} & \S\S\ 13.1, 13.2, 13.3 & \S\S\ 13.1, 13.2, 13.3 \\
\addlinespace[7pt]

31 & \textbf{Du cercle unité à la fonction sinus}\newline {\footnotesize Construire le sinus à partir de la hauteur d’un point sur le cercle.} & \S\S\ 14.2, 16.5 & \S\S\ 14.2, 16.5 \\
\addlinespace[7pt]

32 & \textbf{Trigonométrie du triangle}\newline {\footnotesize Choisir un rapport trigonométrique selon l’angle et les côtés.} & \S\ 15.1 & \S\ 15.1 \\
\addlinespace[7pt]

33 & \textbf{Lois des sinus et des cosinus}\newline {\footnotesize Associer chaque côté à son angle et choisir une loi adaptée.} & \S\S\ 15.2, 15.4 & \S\S\ 15.2, 15.4 \\
\addlinespace[7pt]

34 & \textbf{Fonctions trigonométriques inverses}\newline {\footnotesize Distinguer la valeur principale de l’ensemble des solutions.} & \S\S\ 16.3, 16.4 & \S\S\ 16.3, 16.4 \\
\addlinespace[7pt]

35 & \textbf{Aire du cercle}\newline {\footnotesize Complément géométrique au chapitre 14; la preuve de l’aire du disque dépasse les exercices indiqués.} & Ch. 14 & Ch. 14 \\
\addlinespace[7pt]

36 & \textbf{Pythagore par les aires}\newline {\footnotesize Preuve par les aires en complément des applications aux triangles rectangles.} & Ch. 15 & \S\ 15.1 \\
\addlinespace[7pt]

37 & \textbf{Notation sigma}\newline {\footnotesize Complément transversal sur la lecture et le développement d’une somme; aucune section dédiée dans ces documents.} & Complément transversal & Complément transversal \\
\addlinespace[7pt]

\end{longtable}
\normalsize


\section*{Erreurs fréquentes}


\small
\begin{longtable}{@{}p{0.055\textwidth}p{0.525\textwidth}p{0.17\textwidth}p{0.17\textwidth}@{}}
\rowcolor{soft}\textbf{No} & \textbf{Vidéo et point de révision} & \textbf{Manuel} & \textbf{Cahier} \\
\midrule
\endfirsthead
\rowcolor{soft}\textbf{No} & \textbf{Vidéo et point de révision} & \textbf{Manuel} & \textbf{Cahier} \\
\midrule
\endhead
\midrule
\multicolumn{4}{r}{\footnotesize Suite à la page suivante}\\
\endfoot
\bottomrule
\endlastfoot

38 & \textbf{Implication et équivalence}\newline {\footnotesize Vérifier séparément une implication et sa réciproque.} & \S\ 2.7 & \S\ 2.7 \\
\addlinespace[7pt]

39 & \textbf{Racine d'un produit}\newline {\footnotesize Contrôler les hypothèses dans les réels; la parenthèse sur les complexes dépasse cette section.} & \S\ 1.4 & \S\ 1.6 \\
\addlinespace[7pt]

40 & \textbf{Égalité de fonctions}\newline {\footnotesize Comparer domaine, codomaine et valeurs; l’exercice du cahier cible le domaine après simplification.} & \S\ 3.1 & \S\ 2.2 \\
\addlinespace[7pt]

41 & \textbf{Solutions parasites}\newline {\footnotesize Vérifier les candidats après une mise au carré.} & \S\S\ 2.5, 2.6 & \S\S\ 2.5, 2.6 \\
\addlinespace[7pt]

42 & \textbf{Ordre de composition}\newline {\footnotesize Comparer les deux ordres de composition sur une même entrée.} & \S\ 3.4 & \S\ 3.4 \\
\addlinespace[7pt]

43 & \textbf{Annulation dans les fractions}\newline {\footnotesize Simplifier des facteurs, jamais des termes, et garder les restrictions initiales.} & \S\ 2.2 & \S\ 2.2 \\
\addlinespace[7pt]

\end{longtable}
\normalsize

\end{document}
